\documentclass[dvipdfmx,disablejfam]{jarticle}
\usepackage[dvipdfmx]{graphicx}
\usepackage{url}
\usepackage{amsmath}
\usepackage{amssymb}
\usepackage{bm}
\usepackage{float}
\usepackage{xcolor}
  
% 体裁（必要なら調整）
\topmargin -0.4mm
\oddsidemargin -0.4mm
\evensidemargin -0.4mm
\headheight 0mm
\headsep 0mm
\footskip 10mm
\textheight 247mm
\textwidth 170mm
\setlength{\parindent}{1zw}
\pagestyle{empty}

\makeatletter
\def \fgcaption{\def \@captype{figure} \caption}
\def \tbcaption{\def \@captype{table} \caption}
\makeatother

\def\figurename{Fig.}
\def\tablename{Table}
\def\topfraction{.95}
\def\bottomfraction{.95}
\def\textfraction{.05}
\def\floatpagefraction{.8}
\def\dbltopfraction{.95}
\def\dbltextfraction{.05}
\def\dblfloatpagefraction{.8}

\begin{document}

%==============================
%  ヘッダー
%==============================
\begin{center}
\begin{tabular}{|c|c|c|} \hline
  \parbox{50mm}{\vspace{1mm} \centering{\Large\bf 2026/04/23}} &
  \parbox{55mm}{\vspace{1mm} \centering{\Large\bf M1\quad 市川 詠祥}} &
  \parbox{50mm}{\vspace{1mm} \centering{\Large\bf No.~03}} \\ [2mm] \hline

  \multicolumn{3}{|c|}{
    \parbox{155mm}{
      \centering\Large\bf \vspace{2mm}
       くさび型・円形きずの応力分布
    }
  } \\ [3mm] \hline\hline

  \multicolumn{3}{|c|}{
    \parbox{155mm}{
      \large\vspace{2mm}

      {\Large\bf 今月の目標}
      \begin{itemize}
        \item くさび型きずのT1・T3応力分布の可視化プログラム作成
        \item 円形きず（U字型）のプログラム作成
      \end{itemize}

      {\Large\bf 進捗の要点}
      \begin{itemize}
        \item くさび型きずのT1・T3応力可視化プログラムを作成し, 探触子位置・ピッチ長を変えた比較を実施した
        \item ピッチを1.25\,mmから2.00\,mmに拡大すると,モード変換の様子が明確に確認できた
        \item 探触子位置（端/中心）の変更では応力分布に大きな差は見られなかった
        \item 円形きずの解析プログラムを新規作成中
      \end{itemize}

      {\Large\bf 今後の実施事項}
      \begin{itemize}
        \item 円形きず解析プログラムの完成
        \item 斜め方向応力（$T_\sigma$）の評価
      \end{itemize}

      {\Large\bf 投稿・発表予定}
      \begin{itemize}
        \item 特になし
      \end{itemize}
    }
  } \\ [3mm] \hline
\end{tabular}
\end{center}

\vspace{6mm}
{\Large\bf 今回の内容}
\vspace{3mm}

\subsection*{1. 円形きず解析プログラムの作成}

くさび型のプログラムを参考に, 円形きず（U字型）を対象とした解析プログラムの新規作成をしている途中です.
作成にあたり, ピッチの定義がくさび型と異なる点が問題となっていて,

くさび型ではピッチ1つが三角形（くさび）1個分に対応するのに対し, 
U字型のピッチ定義は$P$ =きず＋隙間＋きず であり, 
1ピッチの中にきず2つ分ある.
このためくさびのプログラムをそのまま流用できず, 形状を書く部分について新たに作成する必要があります．
現在この作業を進めています．

\subsection*{2. くさび型きずの応力分布可視化}

先週に引き続き, くさび型きずを対象としたシミュレーション結果の可視化を行った．

探触子の位置とピッチ長$P$を変化させたときの応力分布（T1: $\sigma_{xx}$, T3: $\sigma_{zz}$）を比較した．

\subsubsection*{2.1 探触子位置の比較（$P = 1.25$\,mm, $d = 0.20$\,mm）}

探触子の中心ラインをくさびの頂点（端）に合わせた場合（Fig.\,\ref{fig:edge125}）と, 
くさびの中心に合わせた場合（Fig.\,\ref{fig:center125}）とを比較した．
両者の応力分布に大きな差は見られなかった．
Fig.\,\ref{fig:points125}に, $P=1.25$\,mm の中心に合わせた場合のT3計測点の配置を示す.

\begin{figure}[H]
  \begin{minipage}[t]{0.49\textwidth}
    \centering
    \includegraphics[width=\textwidth]{../../project4/tmp_output/kusabi_T1_T3_surface_map_edge_pitch125_depth20.png}
    \caption{$P=1.25$\,mm, 探触子を端に配置したときの応力分布}
    \label{fig:edge125}
  \end{minipage}
  \hfill
  \begin{minipage}[t]{0.49\textwidth}
    \centering
    \includegraphics[width=\textwidth]{../../project4/tmp_output/kusabi_T1_T3_surface_map_center_pitch125_depth20.png}
    \caption{$P=1.25$\,mm, 探触子を中心に配置したときの応力分布}
    \label{fig:center125}
  \end{minipage}
\end{figure}

\begin{figure}[H]
  \centering
  \includegraphics[width=0.55\textwidth]{../../project4/tmp_output/kusabi_map_points_center_pitch125_depth20.png}
  \caption{$P=1.25$\,mm, 探触子を中心に配置したときのT3計測点の位置}
  \label{fig:points125}
\end{figure}

\subsubsection*{2.2 ピッチ長の比較（端に配置したとき, $d = 0.20$\,mm）}

ピッチを$P=1.25$\,mmから$P=2.00$\,mmに変更したときの応力分布を比較した
（Fig.\,\ref{fig:edge125}とFig.\,\ref{fig:edge200}）．
ピッチを長くすると, 応力が伝搬する時間が長くなり, $P = 1.25mm $ の
時に比べて, モード変換の様子がわかりやすくなった．

\begin{figure}[H]
  \begin{minipage}[t]{0.49\textwidth}
    \centering
    \includegraphics[width=\textwidth]{../../project4/tmp_output/kusabi_T1_T3_surface_map_edge_pitch125_depth20.png}
    \caption{$P=1.25$\,mm, 探触子を端に配置したときの応力分布（再掲）}
  \end{minipage}
  \hfill
  \begin{minipage}[t]{0.49\textwidth}
    \centering
    \includegraphics[width=\textwidth]{../../project4/tmp_output/kusabi_T1_T3_surface_map_edge_pitch200_depth20.png}
    \caption{$P=2.00$\,mm, 探触子を端に配置したときの応力分布}
    \label{fig:edge200}
  \end{minipage}
\end{figure}

\subsection*{3. 今後の予定}

\begin{itemize}
  \item 円形きず解析プログラムの完成
  \item 斜め方向応力$T_\sigma$の評価
\end{itemize}

%===========================================

\end{document}
